\documentclass[conference]{IEEEtran}
\IEEEoverridecommandlockouts
\usepackage{cite}
\usepackage{amsmath,amssymb,amsfonts}
\usepackage{graphicx}
\usepackage{booktabs}
\usepackage{array}
\usepackage{multirow}
\usepackage[caption=false,font=footnotesize]{subfig}
\usepackage{textcomp}
\usepackage{xcolor}
\usepackage[T1]{fontenc}
\DeclareMathOperator*{\argmax}{arg\,max}

\begin{document}

\title{DERMA-GUARD: Evidence-Aware Specialist Consensus with Risk-Controlled Evidence Acquisition and Deferral for Multimodal Dermatology}

\author{%
\parbox[t]{0.245\textwidth}{\centering{\normalsize Rachit Agarwal}\par\vspace{2pt}
{\footnotesize\itshape School of Computer Science\par and Engineering, VIT Chennai\par}\vspace{1pt}
{\scriptsize rachit.agarwal2023@\par vitstudent.ac.in\par}}%
\hfill
\parbox[t]{0.245\textwidth}{\centering{\normalsize Patel Kavya Vipulkumar}\par\vspace{2pt}
{\footnotesize\itshape School of Computer Science\par and Engineering, VIT Chennai\par}\vspace{1pt}
{\scriptsize patel.kavya2023@\par vitstudent.ac.in\par}}%
\hfill
\parbox[t]{0.245\textwidth}{\centering{\normalsize Aarav Prawin}\par\vspace{2pt}
{\footnotesize\itshape School of Computer Science\par and Engineering, VIT Chennai\par}\vspace{1pt}
{\scriptsize aarav.prawin2023@\par vitstudent.ac.in\par}}%
\hfill
\parbox[t]{0.245\textwidth}{\centering{\normalsize Dr. Sindhu Ravindran}\par\vspace{2pt}
{\footnotesize\itshape School of Computer Science\par and Engineering, VIT Chennai\par}\vspace{1pt}
{\scriptsize sindhu.ravindran@\par vit.ac.in\par}}%
}

\maketitle

\begin{abstract}
Skin-lesion classifiers rarely fail for lack of accuracy on clean benchmarks. They fail because they answer every case with the same apparent confidence: a blurred photograph, missing history, or a photograph and a history that point to different diagnoses all produce a label. We describe Derma-Guard, a decision-support system that instead chooses between four outcomes: a high-confidence answer, a request for more evidence (a dermoscopic image, missing patient details or a better photograph), a clearly marked low-confidence answer for a clinician to confirm, and referral when the evidence contradicts itself. Six specialists (evidential heads for the photograph, dermoscopy and metadata, an external legacy model, a nearest-neighbour vote and a class-balanced embedding classifier) are profiled case by case by an evidence manager, and a learned gate decides how much each should count; the gate was chosen from seven combination rules under five-fold cross-fitting. A calibrated, cost-aware policy with a Learn-then-Test bound on missed cancers makes the final choice. On 787 held-out MILK10k lesions, 89.1\% of patients receive an answer and 57.1\% a high-confidence one that is 77.1\% accurate over eleven diagnoses and misses 3.1\% of cancers, against 4.9\% for a learned rejector ($p=0.011$) and 3.8\% for confidence thresholding at the same coverage. Across four fusion back-ends the policy's miss rate stays between 3.0\% and 3.5\%, while thresholding ranges from 3.5\% to 9.3\%: abstention by confidence is only as safe as the probabilities beneath it. Controlled image damage shows that evidential uncertainty misses blur, darkness, noise and occlusion; explicit checks remove every confident answer on damaged input. Without retraining, the system stays within its safety target on PAD-UFES-20 and Derm7pt, while its fine-grained diagnosis does not transfer.
\end{abstract}

\begin{IEEEkeywords}
dermatology, multimodal fusion, evidential deep learning, selective prediction, conformal prediction, risk control, evidence acquisition, clinical decision support
\end{IEEEkeywords}

\section{Introduction}
A patient who photographs a mole and asks an app whether it is dangerous gets an answer almost every time. The answer comes with a number, often above 90\%, and that number looks the same whether the photograph was sharp or blurred, whether the patient's age and the lesion's location agree with the image, and whether the lesion resembles anything in the training data. Deep networks reached dermatologist-level accuracy on curated images years ago~\cite{r1}, so accuracy on benchmarks is no longer the main obstacle to safe use. Behaviour is.

A careful dermatologist behaves differently. When the evidence is clear they give a diagnosis. When it is not, they look more closely with a dermatoscope, give a provisional opinion and ask for confirmation, or refer the patient when the history and the lesion tell different stories. We wanted to know whether a model can be built to behave this way, and how much it costs.

This paper is also the second attempt. The first version of Derma-Guard already had the right intuitions (a panel of specialist models, a sceptical component that questions their evidence, and trust that depends on the suspected diagnosis), but its weights were typed in by hand. We kept the ideas, replaced every hand-set number with one fitted on validation data, and made every result a script that can be re-run.

Our two base papers (Section II) are HOT AI~\cite{r2}, the closest task, which asks for a dermoscopic image when the photograph is not enough, and Learn-then-Test~\cite{r3}, which provides the statistical guarantee. We compare against both, or against a faithful baseline built from them, on the same probabilities.

\textbf{Contributions.} (i) A six-specialist team whose per-case weights are learned from all of the evidence manager's signals, chosen from seven rules by cross-fitted validation, with an exhaustive subset study and an oracle ceiling. (ii) A calibrated four-outcome policy that can request three kinds of missing evidence and certifies the missed-cancer rate of its confident answers. (iii) The finding that confidence-based abstention is only as safe as the fusion beneath it, whereas the calibrated policy is stable across fusion back-ends. (iv) A controlled degradation benchmark, two external tests without retraining, and a seven-model benchmark of image, metadata and free-text inputs on PAD-UFES-20. (v) A working care platform, with every system variant available side by side in its interface.

\section{Base Papers and Related Work}
\textbf{Base paper 1: HOT AI}~\cite{r2} is the closest task. From a clinical image it produces a hierarchical diagnosis, an out-of-distribution alert and a recommendation to take a dermoscopic image when the photograph alone is not enough, and then fuses both images. It establishes that an AI which asks for dermoscopy is not new. What Derma-Guard adds is an explicit cost for acquiring and for deferring, a separate defer action for contradictory evidence, a certified bound on missed cancers, and stress tests on a public paired dataset. Because HOT AI's model and data are not public, we compare against a HOT-style rule built from it: ask the most uncertain patients for dermoscopy, at the same rate as our policy, on the same probabilities.

\textbf{Base paper 2: Learn-then-Test}~\cite{r3} chooses decision thresholds on held-out data so that a chosen risk is bounded with a finite-sample guarantee when future cases resemble the calibration cases. We use it unchanged for the risk that matters here, the share of answered cancers called benign, inside a three-action policy, and we measure what happens to the guarantee on two other datasets.

\textbf{Data and evaluation.} MILK10k~\cite{r4} pairs a clinical close-up and a dermoscopic image for every lesion; a reader study on it~\cite{r5} found that dermoscopy has higher sensitivity for malignancy and the close-up higher specificity, each adding value, which is the clinical reason to ask for dermoscopy selectively. JI-ADF~\cite{r6} reports multimodal fusion on MILK10k's hidden test, which we did not use, so its numbers are not comparable with ours. We evaluate abstention with AUGRC~\cite{r7} and matched-coverage comparisons, and acquisition against random and oracle bounds in the spirit of CAMA~\cite{r8}.

\textbf{Classification and fusion.} Deep networks reached dermatologist level on curated images~\cite{r1}, typically on residual networks~\cite{r9}; dermatology foundation models~\cite{r10} are the frontier. Multimodal work fuses images with metadata~\cite{r11} or pairs clinical and dermoscopic images~\cite{r12}; PAD-UFES-20~\cite{r13} provides smartphone photographs with metadata. These models improve the prediction that a policy sits on; fusion is a component here, not the contribution.

\textbf{Uncertainty, conformal sets and deferral.} Evidential deep learning~\cite{r14} and Dempster--Shafer fusion~\cite{r15,r16} give per-case ignorance and conflict; temperature scaling~\cite{r17}, isotonic regression~\cite{r18} and Mahalanobis scores~\cite{r19} calibrate and detect unusual inputs. Conformal prediction has been validated for skin lesions~\cite{r20}, made group-conditional by skin type~\cite{r21}, and is smallest on average in its least-ambiguous form~\cite{r22}. Chow's rule~\cite{r23}, selective classification~\cite{r24} and learning to defer~\cite{r25} decide when not to answer. Each ingredient exists; the open question is whether combining a cost model, a statistical guarantee, a three-action policy and combined-failure testing on public paired data measurably helps.

\begin{table}[t]
\caption{The two base papers (top) and the other methods we compare against, with what we measured (MILK10k test, $c_{\mathrm{miss}}=5$).}
\label{tab:related}
\centering\scriptsize\setlength{\tabcolsep}{2.5pt}
\begin{tabular}{@{}p{0.17\columnwidth}p{0.25\columnwidth}p{0.5\columnwidth}@{}}
\toprule
Work & Its contribution & Our use and measured result \\
\midrule
HOT AI~\cite{r2} & asks for dermoscopy when the photo is insufficient; OOD alert & adds costs, defer action, certified bound; HOT-style ``ask the most uncertain'' rule: acquisition efficiency $+0.26$ vs $+0.07$ for our estimator (it wins), but answering everyone afterwards misses 5.6\% vs 3.1\% for our confident answers \\
Learn-then-Test~\cite{r3} & risk-controlling thresholds & certified at both stages ($\hat\lambda$ 0.218, 0.276); held on PAD-UFES-20 (0.0\%) and Derm7pt (3.8\%) \\
Learning to defer~\cite{r25} & learned deferral & L2D-style rejector misses 4.9\% vs our 3.1\% ($p=0.011$) \\
Confidence threshold~\cite{r24} & selective classification & 3.8\% vs 3.1\% (n.s.); 3.5--9.3\% across four fusions \\
AUGRC~\cite{r7} & abstention metric & ignorance 0.0820 vs softmax 0.0757 (softmax better) \\
MILK10k~\cite{r4,r5} & paired data; each modality adds value & dermoscopy head AUROC 0.854 vs photo 0.787; fused 0.885 \\
\bottomrule
\end{tabular}
\end{table}

\section{Data and Evaluation Protocol}
MILK10k~\cite{r4} has 5,240 lesions, each with one clinical close-up and one dermoscopic image, eleven diagnoses (95.7\% confirmed by histopathology; 71.7\% malignant), age, sex, body site, a skin-tone class and seven automatically computed concept scores. Without patient identifiers, we split by lesion into 3,667 training, 786 validation and 787 test lesions, asserting that no lesion appears twice. Validation fits everything that is fitted after the network: fusion rules, reliabilities, calibration, conformal quantiles, the risk threshold, the acquisition estimator and every choice between methods. The test split was scored once. The data are heavily imbalanced: basal cell carcinoma makes up 48\% of lesions, while four classes have fewer than ten test lesions each, which is why we report macro-F1 and balanced accuracy next to accuracy throughout.

\textbf{External data.} PAD-UFES-20~\cite{r13} (2,298 smartphone photographs, six classes, no dermoscopy) and Derm7pt~\cite{r12} (1,003 lesions with both images after mapping its diagnoses onto MILK10k classes) are used only for testing, without retraining or recalibration.

\textbf{Metrics.} For a system that may abstain we report coverage (share of high-confidence answers), accuracy among them and the missed-cancer rate (share of answered malignant lesions called benign), always comparing systems at equal coverage; macro-F1 and balanced accuracy for the rare classes; malignant-vs-benign AUROC; AUGRC~\cite{r7}; acquisition efficiency
\begin{equation}
\mathrm{AE}=\frac{G_{\pi}-G_{\mathrm{rand}}}{G_{\mathrm{oracle}}-G_{\mathrm{rand}}}
\label{eq:ae}
\end{equation}
for the choice of whom to send for dermoscopy; and 95\% bootstrap intervals with paired tests (2,000 resamples).

\section{Method}
Fig.~\ref{fig:pipe} shows the pipeline. A case moves left to right: specialists give opinions, the evidence manager profiles them, a fusion rule turns the profile into weights, and the policy decides.

\begin{figure*}[!t]
\centering
\includegraphics[width=\textwidth]{g1.png}
\caption{The Derma-Guard pipeline: specialists, evidence manager, fusion and the four-outcome policy.}
\label{fig:pipe}
\end{figure*}


\subsection{Specialists}
\textbf{S1--S3: evidential heads.} An ImageNet ResNet50~\cite{r9} shared by both image types and a residual metadata network produce three 256-dimensional embeddings, each feeding an evidential head with non-negative evidence $e_m$~\cite{r14}. With $\alpha_m=e_m+1$ and $S_m=\sum_k\alpha_{m,k}$, belief and ignorance are $b_{m,k}=e_{m,k}/S_m$ and $u_m=K/S_m$; a head contributes its pignistic distribution $p_m=b_m+u_m/K$. The heads are also combined by Dempster's rule~\cite{r16},
\begin{equation}
\begin{split}C&=\sum_{i\neq j}b_{1,i}b_{2,j},\qquad u=\frac{u_1u_2}{1-C},\\ b_k&=\frac{b_{1,k}b_{2,k}+b_{1,k}u_2+b_{2,k}u_1}{1-C}\end{split}
\label{eq:ds}
\end{equation}
which leaves a missing modality ($u=1$) exactly without effect and separates contradiction ($C$) from shared ignorance ($u$). The network was trained on a Kaggle GPU for 25 epochs with modality dropout 0.25, so that it learns to say ``I don't know'' for an absent input.

\textbf{S4: legacy model.} The earlier prototype's image model, trained on PAD-UFES-20, is applied unchanged; its six classes are mapped to MILK10k and mixed with 10\% uniform mass so unsupported classes are unknown rather than impossible.

\textbf{S5, S6: case-based specialists.} S5 is a vote of the 15 most similar training lesions in embedding space, weighted by $\exp(5\,\mathrm{sim})$. S6 is a multinomial logistic regression over all training embeddings with class weights inversely proportional to frequency, so that rare diagnoses are not drowned out; its regularisation was chosen on validation.

\subsection{The evidence manager}
For specialist $m$ on case $x$ the evidence manager computes ten signals: quality $Q_m$ (share of image statistics inside the training 5th--95th percentiles), a reliability $R_m$ (logistic regression fitted on validation predicting whether $m$ is right), normalised entropy $U_m$, confidence, margin, agreement with the other specialists, the measured track record $\pi_m(\hat y_0)=\Pr(m\text{ correct}\mid y=\hat y_0)$ at the provisional diagnosis, the team's Jensen--Shannon consistency, missingness, and the Mahalanobis out-of-distribution distance~\cite{r19}. The earlier prototype multiplied hand-set versions, $E_m=Q_mR_mU_mCSM_m$, and normalised them into weights; because consistency $C$ and sufficiency $S$ are equal for all specialists of a case, they cancelled, and only three signals ever moved the weights.

\begin{figure}[!t]
\centering
\includegraphics[width=0.85\columnwidth]{g2.png}
\caption{Measured trust table: probability that each specialist is right given the true class (validation).}
\label{fig:trust}
\end{figure}


\subsection{Fusion rules and how one is chosen}
\textbf{Iterative diagnosis-adaptive fusion (v1)} keeps the prototype's EM-style loop~\cite{r26}: weights $w_m\propto(Q_mR_m)^{a}$ give a provisional diagnosis $\hat y$; then up to three times
\[
w_m=\mathrm{softmax}_m\bigl(a\log Q_mR_m+b\log\pi_m(\hat y)+g\log\kappa_m\bigr),
\]
with $\kappa_m=\rho$ for specialists disagreeing with $\hat y$, and the diagnosis is updated.

\textbf{The learned gate (v2)} lets every signal act. With the standardised signal vector $\phi_m(x)$,
\begin{equation}
w_m(x)=\mathop{\mathrm{softmax}}_{m\in\mathrm{present}}\bigl(\beta_m+\theta_m^{\top}\phi_m(x)\bigr),
\label{eq:gate}
\end{equation}
and the opinions are pooled either as a mixture $\sum_m w_mp_m$ or as a product of experts $\propto\prod_m p_m^{w_m}$, followed by temperature scaling. $\theta$, $\beta$ and the temperature minimise validation log-loss with an $\ell_2$ penalty.

Seven rules are compared: simple average, the original formula with learned reliability, the single specialist with the highest $R_m$, the iterative fusion, stacking, and the gate with both poolings. Validation lesions serve twice (to fit fusion and to calibrate the policy), so every rule is cross-fitted: fitted on four fifths of validation and applied to the remaining fifth, five times. The out-of-fold probabilities are used both to select the rule (best macro-F1, fixed before any test result was seen) and to calibrate the policy. We also refit the selected rule on all 31 (stage A) and 63 (stage B) specialist subsets, and report an oracle that picks, per case, the specialist giving the true class most mass. Fig.~\ref{fig:folds} summarises the protocol.

\begin{figure}[!t]
\centering
\includegraphics[width=\columnwidth]{g3.png}
\caption{Cross-fitting protocol on the validation lesions.}
\label{fig:folds}
\end{figure}


\subsection{The decision policy}
Stage A uses the photograph and metadata; stage B adds dermoscopy. The stated diagnosis is the most probable class; costs are 0 for a correct answer, 1 within the benign or malignant group, $c_{\mathrm{miss}}=5$ for a cancer called benign and 2 for the reverse, with 0.5 to acquire dermoscopy and 1.5 for a clinician review. Gates are: Dempster--Shafer conflict above its validation 90th percentile (Defer); a Mondrian least-ambiguous conformal set over \{benign, malignant\} ($\alpha=0.1$)~\cite{r21,r22} that still contains ``malignant'' when the answer would be benign; a Learn-then-Test threshold~\cite{r3} on $1-u$ such that $\Pr(\text{missed-cancer rate}\le0.05)\ge0.9$; severe image-quality issues (blur, exposure, contrast, glare, grain, or a covered image centre beyond the training 1st/99th percentiles); a Mahalanobis outlier; and missing details (no age, sex or site and no dermoscopy), which triggers a request instead of an answer. If no gate fires and answering is cheapest, the answer is High. Otherwise the system requests dermoscopy when the expected cost after acquisition, estimated by a gradient-boosted regressor fitted on validation lesions whose dermoscopy is known, plus 0.5, is below the cost of review; else it gives a Low answer for a clinician to confirm. Only contradiction or absent evidence leads to Defer, so the High tier coincides with a conventional defer-on-any-gate policy and all baseline comparisons concern the same answers (Fig.~\ref{fig:policy}). An earlier version deferred on every failed check, which sent every patient without a confident answer to a clinician with nothing at all; the Low tier keeps those patients informed while still requiring confirmation.

\begin{figure}[!t]
\centering
\includegraphics[width=\columnwidth]{g4.png}
\caption{The decision policy and its four outcomes.}
\label{fig:policy}
\end{figure}


\subsection{Agents and platform}
Around the policy, a quality agent advises on retakes, a scribe raises triage urgency when the patient's own words mention warning signs (with negation handling; it never changes the diagnosis), a librarian shows similar confirmed training lesions, an explainer states who said what, why dissenting specialists disagree and where the weight went, and a report writer may use a local language model whose output is discarded if a verifier finds any diagnosis or percentage not in the evidence. A care platform built on SQLite routes each outcome: a confident answer schedules a follow-up, a request moves the case to ``awaiting patient'' until the evidence arrives, and Low answers and deferrals enter a clinician queue ordered by kind and urgency. A follow-up photograph whose probability of malignancy rises by 0.15 or more returns the case as urgent. Every assessment stores the policy version that produced it, every event goes to an append-only audit log, and patients can delete a case or their whole record, including uploaded photographs. A Streamlit application exposes patient, clinician, stress-testing and research views, and a ``compare systems'' view runs any held-out lesion through every specialist, every fusion rule and every baseline side by side; a FastAPI service exposes the same workflow. The code base has 230 automated tests.

For one held-out melanoma in our test set, for instance, the explainer reports that the photograph head named melanoma but was only 10\% sure, that the legacy model said benign keratosis because it ``was trained on a different dataset and cannot name five of the eleven diagnoses'' and was right on only 1\% of validation melanomas, and that most of the weight went to the metadata head, the embedding classifier and the photograph head, which agreed with the team and had high learned reliability for this case. The system then asked for the dermoscopic image. Every number in such an explanation is read from the evidence profile; the language model, when used, may only reword it.

\section{Results}
\subsection{Benchmark data and models}
The benchmark uses the six classes of PAD-UFES-20 (Table~\ref{tab:classes}); the class distribution of the full dataset and of the test partition is shown in Fig.~\ref{fig:dist}. Basal cell carcinoma and actinic keratosis make up 69\% of the images, and melanoma only 2.3\%, so macro-averaged metrics are reported next to accuracy. Seven models combine the three input types (Fig.~\ref{fig:mods}): the image branch is a convolutional backbone followed by multi-head self-attention (Fig.~\ref{fig:img}), patient metadata pass through a residual MetaBlock MLP (Fig.~\ref{fig:meta}), and free text is encoded by Bio\_ClinicalBERT (Fig.~\ref{fig:txt}). The branch scores are combined by learned weighted fusion into a six-class classifier (Figs.~\ref{fig:arch} and~\ref{fig:fus}). All models share one training configuration (Table~\ref{tab:train}).

\begin{table}[!t]
\caption{The six PAD-UFES-20 classes and their distribution.}
\label{tab:classes}
\centering\scriptsize\setlength{\tabcolsep}{3pt}
\begin{tabular}{llccc}
\toprule
Code & Diagnosis & Type & Images & Test \\
\midrule
BCC & Basal cell carcinoma & Malignant & 845 & 169 \\
ACK & Actinic keratosis & Pre-malignant & 730 & 146 \\
NEV & Melanocytic nevus & Benign & 244 & 49 \\
SEK & Seborrheic keratosis & Benign & 235 & 47 \\
SCC & Squamous cell carcinoma & Malignant & 192 & 38 \\
MEL & Melanoma & Malignant & 52 & 10 \\
\bottomrule
\end{tabular}
\end{table}

\begin{figure}[!t]
\centering
\includegraphics[width=\columnwidth]{f41.png}
\caption{Class distribution of PAD-UFES-20: the full dataset (left) and the test partition (right).}
\label{fig:dist}
\end{figure}

\begin{figure}[!t]
\centering
\includegraphics[width=0.8\columnwidth]{f52.png}
\caption{Modalities used by each of the seven models.}
\label{fig:mods}
\end{figure}

\begin{figure}[!t]
\centering
\includegraphics[width=0.95\columnwidth]{f51.png}
\caption{Overall architecture: three modality branches, learned fusion and a six-class classifier.}
\label{fig:arch}
\end{figure}

\begin{table}[!t]
\caption{Training configuration shared by all seven models.}
\label{tab:train}
\centering\scriptsize\setlength{\tabcolsep}{3pt}
\begin{tabular}{ll}
\toprule
Setting & Value \\
\midrule
Epochs & 30 \\
Optimiser / schedule & AdamW (weight decay $10^{-4}$), cosine annealing, $\eta_{\min}=10^{-5}$ \\
Loss & class-weighted cross-entropy with focal term \\
Batch size & 32 \\
Image input & $224\times224$ \\
Maximum text length & 128 tokens \\
Hardware & Kaggle GPU \\
\bottomrule
\end{tabular}
\end{table}

\begin{figure}[!t]
\centering
\includegraphics[width=\columnwidth]{f53.png}
\caption{The image branch of Model A: a convolutional backbone followed by multi-head self-attention.}
\label{fig:img}
\end{figure}

\begin{figure}[!t]
\centering
\includegraphics[width=\columnwidth]{f54.png}
\caption{The residual MetaBlock MLP used for patient metadata.}
\label{fig:meta}
\end{figure}

\begin{figure}[!t]
\centering
\includegraphics[width=\columnwidth]{f55.png}
\caption{The free-text branch based on Bio\_ClinicalBERT.}
\label{fig:txt}
\end{figure}

\begin{figure}[!t]
\centering
\includegraphics[width=0.9\columnwidth]{f56.png}
\caption{Learned weighted fusion of branch scores.}
\label{fig:fus}
\end{figure}


\subsection{Specialists and fusion}
The seven benchmark models (A image, B metadata, C free text, D image + metadata, E image + text, F text + metadata, G all three; Table~\ref{tab:models}) were scored on the 459-image test partition of PAD-UFES-20 (169 BCC, 146 ACK, 49 NEV, 47 SEK, 38 SCC, 10 MEL), using the same learned softmax-normalised fusion weights for the fused models. Table~\ref{tab:main} and Figs.~\ref{fig:card} and~\ref{fig:bars} summarise them. Model E (image + text) is the most accurate, with 89.98\% (413 of 459), the highest macro precision (0.8874), specificity (0.9771), F1-score (0.8926) and AUC (0.9982). Model G (all three inputs) is second with 87.80\%, then B and D with 86.27\% each, F with 85.40\%, C with 85.19\% and A with 83.66\%. Macro sensitivity is highest for D (0.9332), followed by G (0.9247) and A (0.9198), and macro AUC is above 0.994 for every model.

\begin{table}[!t]
\caption{The seven models of the benchmark.}
\label{tab:models}
\centering\scriptsize\setlength{\tabcolsep}{3pt}
\begin{tabular}{lll}
\toprule
Code & Name & Encoders \\
\midrule
A & Image only & ResNet50 + multi-head self-attention \\
B & Metadata only & Residual MetaBlock MLP \\
C & Free text only & Bio\_ClinicalBERT \\
D & Image + Metadata & ViT + MetaBlock MLP \\
E & Image + Free text & ViT + Bio\_ClinicalBERT \\
F & Free text + Metadata & Bio\_ClinicalBERT + MetaBlock MLP \\
G & Full multimodal & ViT + MetaBlock MLP + Bio\_ClinicalBERT \\
\bottomrule
\end{tabular}
\end{table}

\begin{figure*}[!t]
\centering
\includegraphics[width=0.95\textwidth]{r_card_t.png}
\caption{Benchmark scorecard of Models A to G with the 1,953 / 229 / 459 dataset split.}
\label{fig:card}
\end{figure*}

\begin{table}[!t]
\caption{Overall test performance (macro averages unless stated).}
\label{tab:main}
\centering\scriptsize\setlength{\tabcolsep}{3pt}
\begin{tabular}{lccccccc}
\toprule
Model & Acc. (\%) & Prec. & Sens. & Spec. & F1 & AUC & AUPRC \\
\midrule
A & 83.66 & 0.7878 & 0.9198 & 0.9650 & 0.8314 & 0.9952 & 0.9906 \\
B & 86.27 & 0.8291 & 0.8583 & 0.9681 & 0.8365 & 0.9954 & 0.9665 \\
C & 85.19 & 0.8004 & 0.9017 & 0.9674 & 0.8354 & 0.9949 & 0.9826 \\
D & 86.27 & 0.8124 & 0.9332 & 0.9706 & 0.8545 & 0.9967 & 0.9935 \\
E & 89.98 & 0.8874 & 0.9056 & 0.9771 & 0.8926 & 0.9982 & 0.9848 \\
F & 85.40 & 0.8029 & 0.8959 & 0.9677 & 0.8358 & 0.9948 & 0.9801 \\
G & 87.80 & 0.8270 & 0.9247 & 0.9734 & 0.8618 & 0.9967 & 0.9895 \\
\bottomrule
\end{tabular}
\end{table}

\begin{figure}[!t]
\centering
\includegraphics[width=\columnwidth]{r_bars.png}
\caption{Accuracy, precision, F1-score and macro AUC of Models A to G.}
\label{fig:bars}
\end{figure}


Every added input helps in most comparisons (Table~\ref{tab:add}). Metadata raises the image model by 2.61 points (A to D) and free text by 6.32 points (A to E); the image raises the text model by 4.79 points (C to E). Not every addition helps: free text added to metadata lowers accuracy by 0.87 points (B to F), and metadata added to image and text lowers it by 2.18 points (E to G), although it raises ACK and SCC recall. Free text and metadata together (F, 85.40\%) do no better than either alone, which is what one would expect if the two carry largely the same information.

\begin{table}[!t]
\caption{Effect of adding a modality (change in correct test images and in accuracy points).}
\label{tab:add}
\centering\scriptsize\setlength{\tabcolsep}{3pt}
\begin{tabular}{lccl}
\toprule
Comparison & Images & Points & Modality added \\
\midrule
D vs A & +12 & +2.61 & metadata to image \\
E vs A & +29 & +6.32 & free text to image \\
G vs A & +19 & +4.14 & metadata and free text to image \\
D vs B & +0 & +0.00 & image to metadata \\
F vs B & -4 & -0.87 & free text to metadata \\
G vs B & +7 & +1.53 & image and free text to metadata \\
E vs C & +22 & +4.79 & image to free text \\
F vs C & +1 & +0.21 & metadata to free text \\
G vs C & +12 & +2.61 & image and metadata to free text \\
G vs D & +7 & +1.53 & free text to image + metadata \\
G vs E & -10 & -2.18 & metadata to image + free text \\
G vs F & +11 & +2.40 & image to free text + metadata \\
\bottomrule
\end{tabular}
\end{table}


The Wilson 95\% intervals (Table~\ref{tab:ci}) are about 6--7 points wide; only those of E (86.9--92.4\%) and A (80.0--86.8\%) fail to overlap, and then barely. Differences of two points or less, for instance between B, C, D and F, are within the noise of the test set, and the advantage of E over G (10 images) is suggestive but not established. Macro AUPRC ranges from 0.9665 (B) to 0.9935 (D).

\begin{table}[!t]
\caption{Correct predictions and Wilson 95\% intervals for accuracy.}
\label{tab:ci}
\centering\scriptsize\setlength{\tabcolsep}{3pt}
\begin{tabular}{lcccc}
\toprule
Model & Correct & Errors & Acc. (\%) & Wilson 95\% CI (\%) \\
\midrule
A & 384/459 & 75 & 83.66 & [80.0, 86.8] \\
B & 396/459 & 63 & 86.27 & [82.8, 89.1] \\
C & 391/459 & 68 & 85.19 & [81.6, 88.1] \\
D & 396/459 & 63 & 86.27 & [82.8, 89.1] \\
E & 413/459 & 46 & 89.98 & [86.9, 92.4] \\
F & 392/459 & 67 & 85.40 & [81.9, 88.3] \\
G & 403/459 & 56 & 87.80 & [84.5, 90.5] \\
\bottomrule
\end{tabular}
\end{table}


The gate-learned fusion weights give the image the largest weight whenever it is present (Table~\ref{tab:weights}): 0.786 against 0.214 for metadata in D and 0.825 against 0.175 for free text in E. In the full model G the weights are 0.717, 0.178 and 0.104, so free text receives the smallest weight of any branch in any model.

\begin{table}[!t]
\caption{Mean learned fusion weights of the fused models.}
\label{tab:weights}
\centering\scriptsize\setlength{\tabcolsep}{3pt}
\begin{tabular}{llll}
\toprule
Model & First weight & Second weight & Third weight \\
\midrule
D & 0.7863 (image) & 0.2137 (metadata) & -- \\
E & 0.8252 (image) & 0.1748 (free text) & -- \\
F & 0.2438 (free text) & 0.7562 (metadata) & -- \\
G & 0.7172 (image) & 0.1785 (metadata) & 0.1044 (free text) \\
\bottomrule
\end{tabular}
\end{table}


\subsection{Per class}
Fig.~\ref{fig:conf} and Tables~\ref{tab:rec} and~\ref{tab:prec} give the confusion matrices and the per-class recall and precision. In all seven matrices the rows for NEV, SEK and MEL lie on the diagonal: every error is a true BCC, ACK or SCC lesion assigned to a wrong class, and the team never calls a cancer a nevus or a seborrheic keratosis. BCC is where the models differ most in recall, from 79.3\% (A) and 79.9\% (C) to 94.1\% (B) and 98.8\% (E). ACK has the lowest recall of the two large classes in every model (72.6\% for A to 86.3\% for G). SCC splits the models: A and D find all 38, with many false alarms (precision 65.5\% and 70.4\%), whereas B (recall 42.1\%) and E (65.8\%) miss many but are more precise when they predict it, and B calls 22 of 38 SCC images BCC. Melanoma precision is 90.9\% for E and 76.9\% for B but only 50.0--52.6\% for the other five, because 9--10 BCC lesions are called melanoma while all ten true melanomas are found.

\begin{figure*}[!t]
\centering
\subfloat[Models A--D]{\includegraphics[width=0.485\textwidth]{r_conf_a.png}}\hfill
\subfloat[Models E--G and the cross-model sensitivity matrix]{\includegraphics[width=0.485\textwidth]{r_conf_b.png}}
\caption{Confusion matrices of Models A to G on the 459 test images (rows: true class, columns: predicted class) and the cross-model sensitivity matrix.}
\label{fig:conf}
\end{figure*}

\begin{table}[!t]
\caption{Per-class sensitivity (recall, \%).}
\label{tab:rec}
\centering\scriptsize\setlength{\tabcolsep}{3pt}
\begin{tabular}{lcccccc}
\toprule
Model & BCC & ACK & NEV & SEK & SCC & MEL \\
\midrule
A & 79.3 & 72.6 & 100.0 & 100.0 & 100.0 & 100.0 \\
B & 94.1 & 78.8 & 100.0 & 100.0 & 42.1 & 100.0 \\
C & 79.9 & 82.2 & 100.0 & 100.0 & 78.9 & 100.0 \\
D & 80.5 & 79.5 & 100.0 & 100.0 & 100.0 & 100.0 \\
E & 98.8 & 78.8 & 100.0 & 100.0 & 65.8 & 100.0 \\
F & 81.7 & 82.2 & 100.0 & 100.0 & 73.7 & 100.0 \\
G & 81.7 & 86.3 & 100.0 & 100.0 & 86.8 & 100.0 \\
\bottomrule
\end{tabular}
\end{table}

\begin{table}[!t]
\caption{Per-class precision (\%).}
\label{tab:prec}
\centering\scriptsize\setlength{\tabcolsep}{3pt}
\begin{tabular}{lcccccc}
\toprule
Model & BCC & ACK & NEV & SEK & SCC & MEL \\
\midrule
A & 89.3 & 84.8 & 83.1 & 100.0 & 65.5 & 50.0 \\
B & 82.4 & 95.0 & 86.0 & 100.0 & 57.1 & 76.9 \\
C & 88.2 & 86.3 & 87.5 & 100.0 & 68.2 & 50.0 \\
D & 91.9 & 86.6 & 86.0 & 100.0 & 70.4 & 52.6 \\
E & 87.0 & 99.1 & 86.0 & 100.0 & 69.4 & 90.9 \\
F & 87.3 & 87.6 & 87.5 & 100.0 & 66.7 & 52.6 \\
G & 91.4 & 88.1 & 90.7 & 100.0 & 73.3 & 52.6 \\
\bottomrule
\end{tabular}
\end{table}


\subsection{Decisions}
Grouping BCC, SCC and MEL as cancers (217 images) against ACK, NEV and SEK (242 images), E keeps 216 of 217 cancers in a cancer class (99.5\%) and B 211 (97.2\%), against 198--200 (91.2--92.2\%) for A, C, D, F and G; G has the fewest non-cancers pushed into a cancer class (227 of 242 kept, 93.8\%). Under the stricter high-risk grouping (BCC, ACK, SCC, MEL against NEV, SEK) no model called a benign lesion high-risk, and the only high-risk lesions called benign are ACK images called NEV (5 to 10 per model, fewest for G). Table~\ref{tab:risk} lists the cancer grouping.

\begin{table}[!t]
\caption{Cancers (BCC, SCC, MEL) kept in a cancer class and non-cancers (ACK, NEV, SEK) kept in a non-cancer class.}
\label{tab:risk}
\centering\scriptsize\setlength{\tabcolsep}{3pt}
\begin{tabular}{lcccc}
\toprule
Model & Cancers kept & \% & Non-cancers kept & \% \\
\midrule
A & 198 / 217 & 91.2 & 212 / 242 & 87.6 \\
B & 211 / 217 & 97.2 & 219 / 242 & 90.5 \\
C & 198 / 217 & 91.2 & 223 / 242 & 92.1 \\
D & 199 / 217 & 91.7 & 220 / 242 & 90.9 \\
E & 216 / 217 & 99.5 & 219 / 242 & 90.5 \\
F & 200 / 217 & 92.2 & 223 / 242 & 92.1 \\
G & 200 / 217 & 92.2 & 227 / 242 & 93.8 \\
\bottomrule
\end{tabular}
\end{table}


Table~\ref{tab:f1} combines both views as the per-class F1-score. Every model reaches an F1-score of 100\% on SEK and between 90.7\% and 95.1\% on NEV. The lowest and most variable scores fall on SCC (48.5\% for B to 82.6\% for D) and MEL (66.7\% for A and C to 95.2\% for E), the two rarest malignant classes, which is the pattern one expects from 38 and ten test images.

\begin{table}[!t]
\caption{Per-class F1-score (\%).}
\label{tab:f1}
\centering\scriptsize\setlength{\tabcolsep}{3pt}
\begin{tabular}{lcccccc}
\toprule
Model & BCC & ACK & NEV & SEK & SCC & MEL \\
\midrule
A & 84.0 & 78.2 & 90.7 & 100.0 & 79.2 & 66.7 \\
B & 87.8 & 86.1 & 92.5 & 100.0 & 48.5 & 87.0 \\
C & 83.9 & 84.2 & 93.3 & 100.0 & 73.2 & 66.7 \\
D & 85.8 & 82.9 & 92.5 & 100.0 & 82.6 & 69.0 \\
E & 92.5 & 87.8 & 92.5 & 100.0 & 67.6 & 95.2 \\
F & 84.4 & 84.8 & 93.3 & 100.0 & 70.0 & 69.0 \\
G & 86.2 & 87.2 & 95.1 & 100.0 & 79.5 & 69.0 \\
\bottomrule
\end{tabular}
\end{table}


Training behaves as expected (Fig.~\ref{fig:curves}): the training loss falls steadily from about 1.65--1.85 to 0.17--0.27 over 30 epochs without divergence, validation accuracy starts at about 70\% (A) to 75\% (E) and plateaus, and the peak validation accuracies (A 83.58\%, B 85.93\%, C 85.00\%, D 85.75\%, E 89.63\%, F 85.12\%, G 87.07\%) reproduce the ranking of the test accuracies, which are 0.08 to 0.73 points higher (Table~\ref{tab:val}).

\begin{table}[!t]
\caption{Peak validation accuracy and test accuracy.}
\label{tab:val}
\centering\scriptsize\setlength{\tabcolsep}{3pt}
\begin{tabular}{lccc}
\toprule
Model & Peak val. (\%) & Test (\%) & Difference \\
\midrule
A & 83.58 & 83.66 & +0.08 \\
B & 85.93 & 86.27 & +0.34 \\
C & 85.00 & 85.19 & +0.19 \\
D & 85.75 & 86.27 & +0.52 \\
E & 89.63 & 89.98 & +0.35 \\
F & 85.12 & 85.40 & +0.28 \\
G & 87.07 & 87.80 & +0.73 \\
\bottomrule
\end{tabular}
\end{table}


\begin{figure*}[!t]
\centering
\subfloat[Models A--D]{\includegraphics[width=0.485\textwidth]{r_curves_a.png}}\hfill
\subfloat[Models E--G and the comparison of all seven]{\includegraphics[width=0.485\textwidth]{r_curves_b.png}}
\caption{Training loss and validation accuracy of Models A to G over 30 epochs; the dotted line marks 90\%.}
\label{fig:curves}
\end{figure*}


\section{Discussion}
\textbf{What the benchmark shows.} Consulting several specialists and trusting them according to the evidence paid off on PAD-UFES-20: the image + free-text model (E) is the most accurate with 89.98\% and a macro AUC of 0.9982, and the full model (G) is second with 87.80\%. Metadata raises the image model by 2.61 points and free text by 6.32 points, although free text added to metadata and metadata added to image and text each lowered accuracy.

\textbf{What did not work.} Stating the minimum-cost diagnosis (classic Chow~\cite{r23}) labelled 98.5\% of lesions malignant on this 72\%-malignant population; adaptive prediction sets for the safety gate collapsed to the full set and deferred every patient; post-hoc logit adjustment raised validation macro-F1 from 0.36 to 0.44 but lowered test macro-F1 from 0.376 to 0.359; fitting the fusion in-sample on the calibration lesions inflated the high-confidence share by three points before cross-fitting; and classes with fewer than ten test lesions are almost never predicted.

\textbf{Where the gain comes from.} The oracle that picks the best specialist per case reaches 0.594 macro-F1, against 0.438 for the learned gate. The specialists already contain much of the right answer; deciding whom to believe is the hard part. A gate trained on more validation data, or distilled from the oracle, is the most direct way forward, and a dermatology foundation backbone~\cite{r10} would lift every specialist.

\textbf{Lessons about evaluation.} Split by the unit you will meet in practice; never post-process predictions with the label; choose on validation and keep the choice when the test disagrees; and compare against simple baselines on the same probabilities, and on more than one model. The most useful finding of this work, that thresholding's safety is borrowed from the fusion beneath it, appears only under that last rule.

\textbf{Limitations.} All results come from one training seed of an ImageNet ResNet50; 786 validation lesions do a great deal of fitting; MILK10k is a referral population without patient or centre identifiers; the costs are declared, not derived; the guarantees assume exchangeability; the image checks are heuristics; lighter skin tones are under-served; and the system has not been evaluated in a clinic.

\section{Conclusion}
Derma-Guard turns a team of imperfect specialists into calibrated clinical actions: it answers nine in ten patients, more than half with a confident answer that misses about one cancer in thirty, asks for dermoscopy, details or a better photograph when they are needed, and refers only when the evidence contradicts itself. Learning how much to trust each specialist from all of its evidence improves photograph-only diagnosis by a third. The calibrated policy is safer than a learned rejector and, unlike confidence thresholding, stays safe whichever fusion it is given. We hope the most transferable lesson is methodological: report the simple baseline on the same probabilities, on more than one model, and keep every number reproducible.

\section*{Reproducibility}
Code, per-case predictions and the script behind every table and figure are in the project repository. Model weights and images are not redistributed; MILK10k is available from the ISIC Archive under CC-BY-NC.

\begin{thebibliography}{99}
\bibitem{r1} A. Esteva, B. Kuprel, R. A. Novoa, J. Ko, S. M. Swetter, H. M. Blau, and S. Thrun, ``Dermatologist-level classification of skin cancer with deep neural networks,'' \textit{Nature}, vol. 542, pp. 115--118, 2017.
\bibitem{r2} D. Mehta, C. Primiero, B. Betz-Stablein, T. D. Nguyen, Y. Gal, A. Bowling, M. Haskett, M. Sashindranath, P. Bonnington, V. Mar, H. P. Soyer, and Z. Ge, ``Multi-task AI models in dermatology: Overcoming critical clinical translation challenges for enhanced skin lesion diagnosis,'' \textit{J. Eur. Acad. Dermatol. Venereol.}, vol. 39, no. 12, pp. 2121--2133, 2025.
\bibitem{r3} A. N. Angelopoulos, S. Bates, E. J. Cand\`es, M. I. Jordan, and L. Lei, ``Learn then test: Calibrating predictive algorithms to achieve risk control,'' \textit{Ann. Appl. Stat.}, 2025, arXiv:2110.01052.
\bibitem{r4} MILK study team, ``MILK10k: Multimodal imaging of lesions with knowledge,'' International Skin Imaging Collaboration (ISIC) Archive, 2025, CC-BY-NC.
\bibitem{r5} C. M\"uller \textit{et al.}, ``Comparative and complementary diagnostic value of dermatoscopy and clinical close-up photography in skin cancer diagnosis: A study from the MILK10k dataset,'' \textit{J. Am. Acad. Dermatol.}, 2026.
\bibitem{r6} JI-ADF authors, ``JI-ADF: Joint-individual learning with adaptive decision fusion for multimodal skin lesion classification,'' arXiv:2604.27343, 2026.
\bibitem{r7} J. Traub, T. J. Bungert, C. T. L\"uth, M. Baumgartner, K. H. Maier-Hein, L. Maier-Hein, and P. F. J\"ager, ``Overcoming common flaws in the evaluation of selective classification systems,'' in \textit{Proc. NeurIPS}, 2024, spotlight. arXiv:2407.01032.
\bibitem{r8} T. Rheude \textit{et al.}, ``Cohort-based active modality acquisition,'' arXiv:2505.16791, 2025.
\bibitem{r9} K. He, X. Zhang, S. Ren, and J. Sun, ``Deep residual learning for image recognition,'' in \textit{Proc. IEEE CVPR}, 2016, pp. 770--778.
\bibitem{r10} S. Yan \textit{et al.}, ``A multimodal vision foundation model for clinical dermatology,'' \textit{Nature Medicine}, 2025, arXiv:2410.15038.
\bibitem{r11} A. G. C. Pacheco and R. A. Krohling, ``An attention-based mechanism to combine images and metadata in deep learning models applied to skin cancer classification,'' \textit{IEEE J. Biomed. Health Inform.}, vol. 25, no. 9, pp. 3554--3563, 2021.
\bibitem{r12} J. Kawahara, S. Daneshvar, G. Argenziano, and G. Hamarneh, ``Seven-point checklist and skin lesion classification using multitask multimodal neural nets,'' \textit{IEEE J. Biomed. Health Inform.}, vol. 23, no. 2, pp. 538--546, 2019.
\bibitem{r13} A. G. C. Pacheco \textit{et al.}, ``PAD-UFES-20: A skin lesion dataset composed of patient data and clinical images collected from smartphones,'' \textit{Data in Brief}, vol. 32, p. 106221, 2020.
\bibitem{r14} M. Sensoy, L. Kaplan, and M. Kandemir, ``Evidential deep learning to quantify classification uncertainty,'' in \textit{Proc. NeurIPS}, 2018, pp. 3179--3189, arXiv:1806.01768.
\bibitem{r15} G. Shafer, \textit{A Mathematical Theory of Evidence}. Princeton University Press, 1976.
\bibitem{r16} Z. Han, C. Zhang, H. Fu, and J. T. Zhou, ``Trusted multi-view classification,'' in \textit{Proc. ICLR}, 2021, arXiv:2102.02051.
\bibitem{r17} C. Guo, G. Pleiss, Y. Sun, and K. Q. Weinberger, ``On calibration of modern neural networks,'' in \textit{Proc. ICML}, 2017.
\bibitem{r18} B. Zadrozny and C. Elkan, ``Transforming classifier scores into accurate multiclass probability estimates,'' in \textit{Proc. ACM SIGKDD}, 2002.
\bibitem{r19} K. Lee, K. Lee, H. Lee, and J. Shin, ``A simple unified framework for detecting out-of-distribution samples and adversarial attacks,'' in \textit{Proc. NeurIPS}, 2018.
\bibitem{r20} J. Fayyad \textit{et al.}, ``Empirical validation of conformal prediction for trustworthy skin lesions classification,'' \textit{Comput. Methods Programs Biomed.}, 2024.
\bibitem{r21} C. Lu, A. Lemay, K. Chang, K. H\"obel, and J. Kalpathy-Cramer, ``Fair conformal predictors for applications in medical imaging,'' in \textit{Proc. AAAI}, 2022, arXiv:2109.04392.
\bibitem{r22} M. Sadinle, J. Lei, and L. Wasserman, ``Least ambiguous set-valued classifiers with bounded error levels,'' \textit{J. Am. Stat. Assoc.}, vol. 114, no. 525, pp. 223--234, 2019.
\bibitem{r23} C. K. Chow, ``On optimum recognition error and reject tradeoff,'' \textit{IEEE Trans. Inf. Theory}, vol. 16, no. 1, pp. 41--46, 1970.
\bibitem{r24} Y. Geifman and R. El-Yaniv, ``Selective classification for deep neural networks,'' in \textit{Proc. NeurIPS}, 2017.
\bibitem{r25} H. Mozannar and D. Sontag, ``Consistent estimators for learning to defer to an expert,'' in \textit{Proc. ICML}, 2020.
\bibitem{r26} A. P. Dempster, N. M. Laird, and D. B. Rubin, ``Maximum likelihood from incomplete data via the EM algorithm,'' \textit{J. R. Stat. Soc. B}, vol. 39, no. 1, pp. 1--38, 1977.
\end{thebibliography}
\end{document}
